\subsection{All finite nonsingular attachment blocks have a fixed small target}

\begin{claim}{A.26}{Proved}{Transfinite finite-block extensions preserve small detecting images}
Let $\nu$ be an infinite cardinal and let $H$ be a group with
$|H|\leq\nu$. There are a group $E_H$ of cardinality at most
$\nu$ and an embedding $j:H\hookrightarrow E_H$ such that the
following holds for every set-sized ordinal $\theta$.
Let $(H_\alpha)_{\alpha\leq\theta}$ be a continuous tower with
$H_0=H$, group colimits at limit ordinals, and presentations
\[
 H_{\alpha+1}=
 (H_\alpha*F(S_\alpha))/
 \langle\!\langle w_{\alpha,r}:r\in R_\alpha\rangle\!\rangle,
 \qquad \alpha<\theta,
 \tag{A.26.1}
\]
where $S_\alpha$ and $R_\alpha$ are finite, and the exponent
rows of the $w_{\alpha,r}$ in $S_\alpha$ are linearly
independent over $\mathbb Q$. There is a homomorphism
\[
 \Phi:H_\theta\longrightarrow E_H
 \quad\hbox{whose restriction to }H\hbox{ is }j.
 \tag{A.26.2}
\]
The target $E_H$ is independent of the tower and its length.
All the canonical tower maps are injective. Moreover every
nonidentity element of $H_\theta$ survives in a quotient of
cardinality at most $\nu$; these detecting targets may depend
on the chosen element.

In particular, let $Y\subset X$ be an actual two-dimensional
classifying CW pair, with at most $\nu$ cells in $Y$. Suppose
there is a continuous filtration from $Y$ to $X$ in which
each successor adds finitely many one-cells and two-cells,
adds no vertices, and is an ordinary integral homology
isomorphism. Then $\pi_1(X)$ is residually in groups of
cardinality at most $\nu$. This applies to a small homology
core from A.19 whenever the stated finite-block filtration
is supplied; the core construction alone does not supply
such a filtration.
\end{claim}
\begin{proof}
We first construct the target, using the finite nonsingular
coefficient-injectivity theorem
\cite[Theorem~1]{D_OpenAI2026Nonsingular} as in A.19.
Start with $H$. For each of countably many passes, list all
finite nonsingular systems with coefficients in the group
present at the beginning of that pass. There are at most
$\nu$ such systems. Adjoin their universal solution groups
successively, taking group colimits at intervening limit
ordinals. The cited theorem makes every coefficient map
injective. Only finitely many generators are added per
system, so all groups and their final union $E_H$ have
cardinality at most $\nu$. Every finite coefficient list
in the union occurs at the beginning of a later pass.
Consequently every finite nonsingular system over $E_H$
has a solution in $E_H$, and the map $j:H\to E_H$ is
injective.

Construct compatible homomorphisms
$\Phi_\alpha:H_\alpha\to E_H$ by transfinite induction,
starting with $\Phi_0=j$. Given $\Phi_\alpha$, replace
each coefficient in each word $w_{\alpha,r}$ by its
$\Phi_\alpha$-image. This leaves its exponent row in
$S_\alpha$ unchanged. Choose a solution
$(e_s)_{s\in S_\alpha}$ of the resulting finite
nonsingular system in $E_H$. The homomorphism
\[
 H_\alpha*F(S_\alpha)\longrightarrow E_H,
 \qquad h\longmapsto\Phi_\alpha(h),\quad s\longmapsto e_s,
\]
kills every $w_{\alpha,r}$. It therefore factors through
the specified quotient (A.26.1), giving
$\Phi_{\alpha+1}$ with
$\Phi_{\alpha+1}|_{H_\alpha}=\Phi_\alpha$.
At a limit ordinal the universal property of the group
colimit supplies the unique compatible map.
This proves (A.26.2). The same nonsingular theorem proves
injectivity of every successor coefficient map; directed
colimits of these injections preserve all earlier groups.

For the elementwise assertion, use disjoint labels for
the sets $S_\alpha$. Fix, once and for all, a finite word
over $H$ and the earlier labels representing every
coefficient used in (A.26.1). The entire tower then has
the single presentation with coefficient group $H$, all
these variable labels, and all the resulting words as
relators. This identification follows at successors by
the presentation universal property and at limits by
the colimit universal property.

Let $g\ne1$ in $H_\theta$. Choose a finite word for it
and include the indices of the blocks supplying its
variable letters. Close this set of indices under the
following operation: for an included block, include
every earlier block whose variable letters occur in
the chosen coefficient words of that block. Each
operation adds finitely many indices per included
block. After countably many iterations the resulting
set $J\subset\theta$ is countable and closed under
these dependencies.

Let $K$ be the presentation over $H$ using precisely
the blocks indexed by $J$, in their inherited order.
Their coefficient words use only $H$ and earlier
blocks in $J$, so this is a valid tower of the form
(A.26.1). Its cardinality is at most $\nu$. Reorder
the full presentation by first attaching these blocks
and then attaching the blocks in $\theta\setminus J$
in their original order. An unselected block has all
its coefficient letters available: selected letters
are already present, and its unselected dependencies
precede it. Its exponent rows in its own new variables
remain unchanged. The reordered presentation has the
same generators and the same relators as the original
one, and hence the identity on these labels gives its
canonical isomorphism to $H_\theta$. Nonsingular
coefficient injection at every remaining step proves
that $K\to H_\theta$ is injective.

The chosen word for $g$ belongs to $K$ and represents
a nonidentity element there. Apply (A.26.2), now with
initial group $K$, to this reordered remaining tower.
It gives $H_\theta\to E_K$ injective on $K$, with
$|E_K|\leq\nu$. Restricting the codomain to the image
gives a quotient of the required size detecting $g$.
The existence of this varying target does not assert
that the original fixed map to $E_H$ detects every
element.

Finally, in the stated CW filtration, van Kampen
gives exactly the presentations (A.26.1). The relative
cellular chain complex at a successor is a finite
free group in degree two mapping to a finite free
group in degree one. Its homology vanishes, so that
boundary is an isomorphism. Its columns are the
relative exponent vectors, which are therefore
independent. Van Kampen also identifies the
fundamental group of each actual union with the
corresponding group colimit. Since the number of
cells of $Y$ bounds $|\pi_1(Y)|$ by $\nu$, the
algebraic assertion applies and proves the stated
residual bound.
\end{proof}
